\documentclass{amsart}
% For dealing with references we use the comment environment
\usepackage{verbatim}
\newenvironment{reference}{\comment}{\endcomment}
%\newenvironment{reference}{}{}
\newenvironment{slogan}{\comment}{\endcomment}
\newenvironment{history}{\comment}{\endcomment}

% For commutative diagrams we use Xy-pic
\usepackage[all]{xy}

% We use 2cell for 2-commutative diagrams.
\xyoption{2cell}
\UseAllTwocells

% We use multicol for the list of chapters between chapters
\usepackage{multicol}

% This is generally recommended for better output
\usepackage{lmodern}
\usepackage[T1]{fontenc}

% For cross-file-references

% Package for hypertext links:
\usepackage{hyperref}

% For any local file, say "hello.tex" you want to link to please

% Theorem environments.
%
\theoremstyle{plain}
\newtheorem{theorem}[subsection]{Theorem}
\newtheorem{proposition}[subsection]{Proposition}
\newtheorem{lemma}[subsection]{Lemma}

\theoremstyle{definition}
\newtheorem{definition}[subsection]{Definition}
\newtheorem{example}[subsection]{Example}
\newtheorem{exercise}[subsection]{Exercise}
\newtheorem{situation}[subsection]{Situation}

\theoremstyle{remark}
\newtheorem{remark}[subsection]{Remark}
\newtheorem{remarks}[subsection]{Remarks}

\numberwithin{equation}{subsection}

% Macros
%
\def\lim{\mathop{\mathrm{lim}}\nolimits}
\def\colim{\mathop{\mathrm{colim}}\nolimits}
\def\Spec{\mathop{\mathrm{Spec}}}
\def\Hom{\mathop{\mathrm{Hom}}\nolimits}
\def\Ext{\mathop{\mathrm{Ext}}\nolimits}
\def\SheafHom{\mathop{\mathcal{H}\!\mathit{om}}\nolimits}
\def\SheafExt{\mathop{\mathcal{E}\!\mathit{xt}}\nolimits}
\def\Sch{\mathit{Sch}}
\def\Mor{\mathop{\mathrm{Mor}}\nolimits}
\def\Ob{\mathop{\mathrm{Ob}}\nolimits}
\def\Sh{\mathop{\mathit{Sh}}\nolimits}
\def\NL{\mathop{N\!L}\nolimits}
\def\CH{\mathop{\mathrm{CH}}\nolimits}
\def\proetale{{pro\text{-}\acute{e}tale}}
\def\etale{{\acute{e}tale}}
\def\QCoh{\mathit{QCoh}}
\def\Ker{\mathop{\mathrm{Ker}}}
\def\Im{\mathop{\mathrm{Im}}}
\def\Coker{\mathop{\mathrm{Coker}}}
\def\Coim{\mathop{\mathrm{Coim}}}

% Boxtimes
%
\DeclareMathSymbol{\boxtimes}{\mathbin}{AMSa}{"02}

%
% Macros for moduli stacks/spaces
%
\def\QCohstack{\mathcal{QC}\!\mathit{oh}}
\def\Cohstack{\mathcal{C}\!\mathit{oh}}
\def\Spacesstack{\mathcal{S}\!\mathit{paces}}
\def\Quotfunctor{\mathrm{Quot}}
\def\Hilbfunctor{\mathrm{Hilb}}
\def\Curvesstack{\mathcal{C}\!\mathit{urves}}
\def\Polarizedstack{\mathcal{P}\!\mathit{olarized}}
\def\Complexesstack{\mathcal{C}\!\mathit{omplexes}}
% \Pic is the operator that assigns to X its picard group, usage \Pic(X)
% \Picardstack_{X/B} denotes the Picard stack of X over B
% \Picardfunctor_{X/B} denotes the Picard functor of X over B
\def\Pic{\mathop{\mathrm{Pic}}\nolimits}
\def\Picardstack{\mathcal{P}\!\mathit{ic}}
\def\Picardfunctor{\mathrm{Pic}}
\def\Deformationcategory{\mathcal{D}\!\mathit{ef}}

\setlength{\emergencystretch}{3em}
\begin{document}
\title[Stacks Project, Tag 0E2X]{261 --- Global existence of relative dualizing complexes for flat locally finitely presented morphisms}
\maketitle
\noindent Copyright (C) 2005--2025 Johan de Jong. Stacks Project authors. Source: \href{https://stacks.math.columbia.edu/tag/0E2X}{Tag 0E2X}.

\noindent Editorial difficulty note (not author text). Difficulty H2: Local algebraic relative dualizing objects are glued using uniqueness and vanishing of negative self-Ext. The argument must then separately construct the global rigidifying map on a neighbourhood of the diagonal by gluing degree-zero derived-Hom sections. Producing the object alone does not satisfy the theorem..

\noindent Editorial provenance note (not author text). History: excerpt from revision \texttt{a0444\allowbreak 6e57e\allowbreak c1fbc\allowbreak 252a8\allowbreak 71afc\allowbreak ec775\allowbreak 2fb28\allowbreak 07b14}, duality.tex, lines 7161--7251. Reference presentation was adapted; original-author context and core support blocks were retained or added. The mathematical wording is unchanged. Licensed under GNU FDL 1.2 or later; no invariant sections or cover texts. See ../licenses/COPYING. Context blocks reproduce author definitions and supporting results; they are not additional counted problems.

\begin{definition}
\label{definition-relative-dualizing-complex}
Let $X \to S$ be a morphism of schemes which is flat and
locally of finite presentation. Let $W \subset X \times_S X$
be any open such that the diagonal $\Delta_{X/S} : X \to X \times_S X$
factors through a closed immersion $\Delta : X \to W$.
A {\it relative dualizing complex} is a
pair $(K, \xi)$ consisting of an object $K \in D(\mathcal{O}_X)$
and a map
$$
\xi : \Delta_*\mathcal{O}_X \longrightarrow L\text{pr}_1^*K|_W
$$
in $D(\mathcal{O}_W)$ such that
\begin{enumerate}
\item $K$ is $S$-perfect (Derived Categories of Schemes, Definition
\ref{perfect-definition-relatively-perfect}), and
\item $\xi$ defines an isomorphism of $\Delta_*\mathcal{O}_X$
with
$R\SheafHom_{\mathcal{O}_W}(
\Delta_*\mathcal{O}_X, L\text{pr}_1^*K|_W)$.
\end{enumerate}
\end{definition}

\begin{definition}
\label{perfect-definition-relatively-perfect}
Let $f : X \to S$ be a morphism of schemes which is flat and
locally of finite presentation. An object $E$ of $D(\mathcal{O}_X)$ is
{\it perfect relative to $S$} or
{\it $S$-perfect} if $E$ is pseudo-coherent
(Cohomology, Definition \ref{cohomology-definition-pseudo-coherent}) and
$E$ locally has finite tor dimension as an object of
$D(f^{-1}\mathcal{O}_S)$
(Cohomology, Definition \ref{cohomology-definition-tor-amplitude}).
\end{definition}

\begin{definition}
\label{cohomology-definition-pseudo-coherent}
Let $(X, \mathcal{O}_X)$ be a ringed space. Let $\mathcal{E}^\bullet$
be a complex of $\mathcal{O}_X$-modules. Let $m \in \mathbf{Z}$.
\begin{enumerate}
\item We say $\mathcal{E}^\bullet$ is {\it $m$-pseudo-coherent}
if there exists an open covering $X = \bigcup U_i$ and for each $i$
a morphism of complexes
$\alpha_i : \mathcal{E}_i^\bullet \to \mathcal{E}^\bullet|_{U_i}$
where $\mathcal{E}_i^\bullet$ is strictly perfect on $U_i$ and
$H^j(\alpha_i)$ is an isomorphism for $j > m$ and $H^m(\alpha_i)$
is surjective.
\item We say $\mathcal{E}^\bullet$ is {\it pseudo-coherent}
if it is $m$-pseudo-coherent for all $m$.
\item We say an object $E$ of $D(\mathcal{O}_X)$ is
{\it $m$-pseudo-coherent} (resp.\ {\it pseudo-coherent})
if and only if it can be represented by a $m$-pseudo-coherent
(resp.\ pseudo-coherent) complex of $\mathcal{O}_X$-modules.
\end{enumerate}
\end{definition}

\begin{definition}
\label{cohomology-definition-tor-amplitude}
Let $(X, \mathcal{O}_X)$ be a ringed space.
Let $E$ be an object of $D(\mathcal{O}_X)$.
Let $a, b \in \mathbf{Z}$ with $a \leq b$.
\begin{enumerate}
\item We say $E$ has {\it tor-amplitude in $[a, b]$}
if $H^i(E \otimes_{\mathcal{O}_X}^\mathbf{L} \mathcal{F}) = 0$
for all $\mathcal{O}_X$-modules $\mathcal{F}$ and all $i \not \in [a, b]$.
\item We say $E$ has {\it finite tor dimension}
if it has tor-amplitude in $[a, b]$ for some $a, b$.
\item We say $E$ {\it locally has finite tor dimension}
if there exists an open covering $X = \bigcup U_i$ such that
$E|_{U_i}$ has finite tor dimension for all $i$.
\end{enumerate}
An $\mathcal{O}_X$-module $\mathcal{F}$ has {\it tor dimension $\leq d$}
if $\mathcal{F}[0]$ viewed as an object of $D(\mathcal{O}_X)$ has
tor-amplitude in $[-d, 0]$.
\end{definition}

\begin{lemma}
\label{lemma-existence-relative-dualizing}
Let $X \to S$ be a morphism of schemes which is
flat and locally of finite presentation.
There exists a relative dualizing complex $(K, \xi)$.
\end{lemma}

\begin{proof}
Let $\mathcal{B}$ be the collection of affine opens of
$X$ which map into an affine open of $S$. For each $U$
we have a relative dualizing complex $(K_U, \xi_U)$ for
$U$ over $S$. Namely, choose an affine open
$V \subset S$ such that $U \to X \to S$ factors through $V$.
Write $U = \Spec(A)$ and $V = \Spec(R)$. By
Dualizing Complexes, Lemma \href{https://stacks.math.columbia.edu/tag/0E2F}{lemma base change relative dualizing (Tag 0E2F)}
there exists a relative dualizing complex $K_A \in D(A)$
for $R \to A$. Arguing backwards through the proof of
Lemma \href{https://stacks.math.columbia.edu/tag/0E2U}{lemma relative dualizing complex algebra (Tag 0E2U)}
this determines an $V$-perfect object $K_U \in D(\mathcal{O}_U)$
and a map
$$
\xi : \Delta_*\mathcal{O}_U \to L\text{pr}_1^*K_U
$$
in $D(\mathcal{O}_{U \times_V U})$. Since being $V$-perfect is the
same as being $S$-perfect and since $U \times_V U = U \times_S U$
we find that $(K_U, \xi_U)$ is as desired.

\medskip\noindent
If $U' \subset U \subset X$ with $U', U \in \mathcal{B}$, then
we have a unique isomorphism $\rho_{U'}^U : K_U|_{U'} \to K_{U'}$
in $D(\mathcal{O}_{U'})$ sending $\xi_U|_{U' \times_S U'}$
to $\xi_{U'}$ by Lemma \href{https://stacks.math.columbia.edu/tag/0E2W}{lemma uniqueness relative dualizing (Tag 0E2W)}
(note that trivially the restriction of a relative dualizing
complex to an open is a relative dualizing complex).
The uniqueness guarantees that
$\rho^U_{U''} = \rho^V_{U''} \circ \rho ^U_{U'}|_{U''}$
for $U'' \subset U' \subset U$ in $\mathcal{B}$.
Observe that $\text{Ext}^i(K_U, K_U) = 0$ for $i < 0$
for $U \in \mathcal{B}$ by Lemma \href{https://stacks.math.columbia.edu/tag/0E2V}{lemma relative dualizing RHom (Tag 0E2V)}
applied to $U/S$ and $K_U$.
Thus the BBD glueing lemma
(Cohomology, Theorem \href{https://stacks.math.columbia.edu/tag/0D6C}{theorem glueing bbd general (Tag 0D6C)})
tells us there is a unique solution, namely, an object
$K \in D(\mathcal{O}_X)$ and isomorphisms $\rho_U : K|_U \to K_U$
such that we have
$\rho^U_{U'} \circ \rho_U|_{U'} = \rho_{U'}$ for all $U' \subset U$,
$U, U' \in \mathcal{B}$.

\medskip\noindent
To finish the proof we have to construct the map
$$
\xi : \Delta_*\mathcal{O}_X \longrightarrow L\text{pr}_1^*K|_W
$$
in $D(\mathcal{O}_W)$ inducing an isomorphism from $\Delta_*\mathcal{O}_X$ to
$R\SheafHom_{\mathcal{O}_W}(\Delta_*\mathcal{O}_X, L\text{pr}_1^*K|_W)$.
Since we may change $W$, we choose
$W = \bigcup_{U \in \mathcal{B}} U \times_S U$.
We can use $\rho_U$ to get isomorphisms
$$
R\SheafHom_{\mathcal{O}_W}(
\Delta_*\mathcal{O}_X, L\text{pr}_1^*K|_W)|_{U \times_S U}
\xrightarrow{\rho_U}
R\SheafHom_{\mathcal{O}_{U \times_S U}}(
\Delta_*\mathcal{O}_U, L\text{pr}_1^*K_U)
$$
As $W$ is covered by the opens $U \times_S U$
we conclude that the cohomology sheaves of
$R\SheafHom_{\mathcal{O}_W}(\Delta_*\mathcal{O}_X, L\text{pr}_1^*K|_W)$
are zero except in degree $0$. Moreover, we obtain isomorphisms
$$
H^0\left(U \times_S U, R\SheafHom_{\mathcal{O}_W}(\Delta_*\mathcal{O}_X,
L\text{pr}_1^*K|_W)\right)
\xrightarrow{\rho_U}
H^0\left((R\SheafHom_{\mathcal{O}_{U \times_S U}}(
\Delta_*\mathcal{O}_U, L\text{pr}_1^*K_U)\right)
$$
Let $\tau_U$ in the LHS be an element mapping to $\xi_U$ under this map.
The compatibilities between
$\rho^U_{U'}$, $\xi_U$, $\xi_{U'}$, $\rho_U$, and $\rho_{U'}$
for $U' \subset U \subset X$ open $U', U \in \mathcal{B}$
imply that $\tau_U|_{U' \times_S U'} = \tau_{U'}$.
Thus we get a global section $\tau$ of the $0$th cohomology sheaf
$H^0(R\SheafHom_{\mathcal{O}_W}(\Delta_*\mathcal{O}_X, L\text{pr}_1^*K|_W))$.
Since the other cohomology sheaves of
$R\SheafHom_{\mathcal{O}_W}(\Delta_*\mathcal{O}_X, L\text{pr}_1^*K|_W)$
are zero, this global section $\tau$
determines a morphism $\xi$ as desired. Since the restriction
of $\xi$ to $U \times_S U$ gives $\xi_U$, we see that it
satisfies the final condition of
Definition \ref{definition-relative-dualizing-complex}.
\end{proof}
\end{document}
